%%--------------------------------------------------------------------
%1--------------------------------------------------------------------
\vspace*{\fill}
\begin{glyph}{Near (近)}{near}\label{glyph:near}
Near.\\\hline
    The sense of \textit{``recent''} is also incorporated.
\end{glyph}

\begin{comps}
\comprtwocone & 斤 + 辵辶⻌⻍ \\\hline
    \comprthreecone & Axe + Walk/Road  \\\hline
\end{comps}

\begin{remglyph}
    \hooglig{Nearly} got \remcom{axed} while \remcom{walking}.\\\hline
\end{remglyph}

\begin{prons}
\pronsrtwocone & \textbf{キン (KIN)} & \textbf{ちか (chika)}\\\hline
\pronsrthreecone & --- & --- \\\hline
\end{prons}

\begin{remprons}
    \hooglig{Near} \comon{kindred} to us are the \comkun{chickadees}.\\\hline
\end{remprons}

%{\middel}
% &  ()\\\hline
\begin{somme}
    近い  &  \textit{near}   &  ちか{\middel}い \mbox{(chika{\middel}i)}\\\hline
    近づく  &  \textit{to approach} &  ちか{\middel}ずく \mbox{(chika{\middel}zuku)}\\\hline
    近道  &  near + road = \textit{shortcut} &  ちか{\middel}みち \mbox{(chika{\middel}michi)} \\\hline
    近親  &  near + parent = \textit{close relative} &  キン{\middel}シン \mbox{(KIN{\middel}SHIN)} \\\hline
    \notyet{最近}  &  most + near = \textit{recently} &  サイ{\middel}キン \mbox{(SAI{\middel}KIN)} \\\hline
\end{somme}

\begin{sentence}
    \underline{\ruby{新}{あたら}しい}{\;}\underline{\ruby{家}{いえ}}\ruby{は}{わ}\underline{\ruby{川}{かわ}}の\underline{\ruby{近}{ちか}く}に\underline{あります}。\\\hline
atara{\middel}shii ie wa kawa no chika{\middel}ku ni arimasu.\\\hline
    (new) (house) (river) (near) (is)\\\hline
    =\textit{The new house is near a river.}\\\hline
\end{sentence}
\end{multicols}
\vhrulefill{5pt}
%%--------------------------------------------------------------------
%2--------------------------------------------------------------------
\begin{multicols}{2}
\begin{glyph}{Lose (失)}{lose}\label{glyph:lose}
Lose, in the sense of ``mislay''.\\\hline
This character often conveys a hint of ``failure and error''.
\end{glyph}

\begin{comps}
\comprtwocone & 丿 + 二 + 人  \\\hline
    \comprthreecone & Slash + Two + Person\\\hline
\end{comps}

\begin{remglyph}
    \hooglig{Lost} \remcom{two} \remcom{persons} \remcom{slashing} each other.\\\hline
\end{remglyph}

\begin{prons}
\pronsrtwocone & \textbf{シツ (SHITSU)} & \textbf{うしな (ushina)}\\\hline
\pronsrthreecone & --- & --- \\\hline
\end{prons}

\begin{remprons}
    \hooglig{Lose} your \comon{sheets} until \comkun{you shine}.\\\hline
\end{remprons}

%{\middel}
% &  ()\\\hline
\begin{somme}
    失う  & \textit{to lose}  &  うしな{\middel}う \mbox{(ushina{\middel}u)}\\\hline
    見失う  &  see + lose = \textit{to lose sight of} & み　うしな{\middel}う \mbox{(mi ushina{\middel}u)} \\\hline
    失言  &  lose + say = \textit{slip of the tongue} & シス{\middel}ゲン \mbox{(SHITSU{\middel}GEN)} \\\hline
\end{somme}

\begin{sentence}
    \underline{\ruby{大会}{タイカイ}}で\underline{\ruby{失言}{シツゲン}}しました。\\\hline
TAI{\middel}KAI de SHITSU{\middel}GEN shimashita.\\\hline
    (convention) (slip of the tongue) (did)\\\hline
    =\textit{(I) made a slip of the tongue at the convention.} \\\hline
\end{sentence}
\end{multicols}
\vspace*{\fill}
\pagebreak
%%--------------------------------------------------------------------
%3--------------------------------------------------------------------
\vspace*{\fill}
\begin{multicols}{2}
\begin{glyph}{Bow (弓)}{bow}\label{glyph:bow}
    Bow (the partner to an arrow).
\end{glyph}

\begin{comps}
\comprtwocone & 弓 \\\hline
    \comprthreecone & Bow \\\hline
\end{comps}

\begin{remglyph}
Part of the Kanji radicals (\#{57}).\\\hline
\end{remglyph}

\begin{prons}
    \pronsrtwocone & \textbf{キュウ (KY\={U})} & \textbf{ゆみ (yumi)}\\\hline
\pronsrthreecone & --- & --- \\\hline
\end{prons}

\begin{remprons}
    This \hooglig{bow} will \comon{kill you} \comkun{yummily}.\\\hline
\end{remprons}

%{\middel}
% &  ()\\\hline
\begin{somme}
    弓  & \textit{bow}  &  ゆみ \mbox{(yumi)} \\\hline
    弓道  & bow + road = \textit{(Japanese) archery}  & キュウ{\middel}ドウ \mbox{(KY\={U}{\middel}D\={O})}  \\\hline
\end{somme}

\begin{sentence}
    \underline{\ruby{前田}{まえだ}さん}\ruby{は}{わ}\underline{\ruby{古}{ふる}い}{\;}\underline{\ruby{弓}{ゆみ}}と\underline{\ruby{美}{うつく}しい}\underline{\ruby{日本刀}{ニホントウ}}\ruby{を}{お}\underline{\ruby{買}{か}いました}。\\\hline
    Mae{\middel}da-san wa furu{\middel}i yumi to utsuku{\middel}shii NI{\middel}HON{\middel}T\={O} o ka{\middel}imashita.\\\hline
    (Maeda-san) (old) (bow) (beautiful) (Japanese sword) (bought)\\\hline
    =\textit{Maeda-san bough an old bow and a beautiful Japanese sword.}\\\hline
\end{sentence}
\end{multicols}
\vhrulefill{5pt}
%%--------------------------------------------------------------------
%4--------------------------------------------------------------------
\begin{multicols}{2}
\begin{glyph}{Horse (馬)}{horse}\label{glyph:horse}
Horse.
\end{glyph}

\begin{comps}
\comprtwocone & 馬 \\\hline
    \comprthreecone & Horse \\\hline
\end{comps}

\begin{remglyph}
Part of the Kanji radicals (\#{187}).\\\hline
\end{remglyph}

\begin{prons}
\pronsrtwocone & \textbf{バ (BA)} & \textbf{うま (uma)}\\\hline
\pronsrthreecone & --- & --- \\\hline
\end{prons}

\begin{remprons}
    \hooglig{Horse}-\ucomon{bum} \ucomkun{mud} is \comkun{umami}.\\\hline
\end{remprons}

%{\middel}
% &  ()\\\hline
\begin{somme}
    馬  & \textit{horse}  &  うま \mbox{(uma)} \\\hline
    子馬  & child + horse = \textit{colt; pony}  & こ{\middel}うま \mbox{(ko{\middel}uma)}  \\\hline
    牛馬  & cow + horse = \textit{cows and horses}  & ギュウ{\middel}バ \mbox{(GY\={U}{\middel}BA)}  \\\hline
    \notyet{馬具}  & horse + tool = \textit{harness}  &  バ{\middel}グ \mbox{(BA{\middel}GU)} \\\hline
\end{somme}

\begin{sentence}
    \underline{\ruby{馬}{うま}}と\underline{\ruby{牛}{うし}}の\underline{どちら}が\underline{\ruby{好}{す}き}ですか。\\\hline
uma to ushi no dochira ga su{\middel}ki desu ka.\\\hline
    (horse) (cow) (which) (like)\\\hline
    =\textit{Which do you like, horses or cows?}\\\hline
\end{sentence}
\end{multicols}
\vspace*{\fill}
\pagebreak
%%--------------------------------------------------------------------
%5--------------------------------------------------------------------
\vspace*{\fill}
\begin{multicols}{2}
\begin{glyph}{Week (週)}{week}\label{glyph:week}
Week.
\end{glyph}

\begin{comps}
\comprtwocone & 周 + 辵辶⻌ \\\hline
    \comprthreecone & Around (\ref{glyph:around}) + Walk/Road \\\hline
\end{comps}

\begin{remglyph}
    \hooglig{Weeks} of \remcom{walking} \remcom{around}.\\\hline
\end{remglyph}

\begin{prons}
    \pronsrtwocone & \textbf{シュウ (SH\={U})} & ---\\\hline
\pronsrthreecone & --- & --- \\\hline
\end{prons}

\begin{remprons}
    \hooglig{Weekly} \comon{shoes}.\\\hline
\end{remprons}

%{\middel}
% &  ()\\\hline
\begin{somme}
    週間  & week + interval = \textit{week}  & シュウ{\middel}カン \mbox{(SH\={U}{\middel}KAN)}  \\\hline
    一週間  & one + week + interval = \textit{one week}  & イッ{\middel}シュウ{\middel}カン \mbox{(IS{\middel}SH\={U}{\middel}KAN)}  \\\hline
    二週間  & two + week + interval = \textit{two weeks}  & ニ{\middel}シュウ{\middel}カン \mbox{(NI{\middel}SH\={U}{\middel}KAN)}  \\\hline
    週日  & week + sun = \textit{weekday}  & シュウ{\middel}ジツ \mbox{(SH\={U}{\middel}JITSU)}  \\\hline
    先週  & precede + week = \textit{last week}  & セン{\middel}シュウ \mbox{(SEN{\middel}SH\={U})}  \\\hline
    \notyet{来週}  & come + week = \textit{next week}  & ライ{\middel}シュウ \mbox{(RAI{\middel}SH\={U})}  \\\hline
\end{somme}

\begin{sentence}
    \underline{\ruby{月曜日}{ゲツヨウび}}{\;}\underline{から}{\;}\underline{\ruby{二週間}{ニシュウカン}}の\underline{\ruby{夏休}{なつやす}み}です。\\\hline
    GETSU{\middel}Y\={O}{\middel}bi kara NI{\middel}SH\={U}{\middel}KAN no natsu yasu{\middel}mi desu.\\\hline
    (Monday) (from) (two weeks) (summer vacation)\\\hline
    =\textit{There's a two-week summer vacation from Monday.}\\\hline
\end{sentence}
\end{multicols}
\vhrulefill{5pt}
%%--------------------------------------------------------------------
%6--------------------------------------------------------------------
\begin{multicols}{2}
\begin{glyph}{Pull (引)}{pull}\label{glyph:pull}
Pull.\\\hline
    This character also conveys the idea of ``attraction'' (best seen in the
    fourth example), and the sense of ``receding'' and ``withdrawing''.\\\hline
    You will notice it frequently on signs; it indicates ``pull'' on doors, and
    ``discount'' when used with numbers.
\end{glyph}

\begin{comps}
\comprtwocone & 弓 + 丨 \\\hline
    \comprthreecone & Bow + Line/Pole/Vertical stroke \\\hline
\end{comps}

\begin{remglyph}
    \textit{Somewhat obvious from the glyph itself.}\\\hline
\end{remglyph}

\begin{prons}
\pronsrtwocone & \textbf{イン (IN)} & \textbf{ひ (hi)}\\\hline
\pronsrthreecone & --- &  ---\\\hline
\end{prons}

\begin{remprons}
    \hooglig{Pull} \comkun{heaps} of \comon{ink}.\\\hline
\end{remprons}

%{\middel}
% &  ()\\\hline
\begin{somme}
    \rye{3}{As seen in the third example, this is another kanji that can act
    like 買.  The 訓読み can also become voiced in the second position, as
    happens in the fifth compound.}
    引く  & \textit{to pull}  & ひ{\middel}く \mbox{(hi{\middel}ku)}  \\\hline
    引き出す  & pull + exit = \textit{to withdraw (something)}  & ひ{\middel}き　だ{\middel}す \mbox{(hi{\middel}ki da{\middel}su)}  \\\hline
    引出し  & pull + exit = \textit{a withdrawal; a drawer}  & ひ{\middel}き　だ{\middel}し \mbox{(hi{\middel}ki da{\middel}shi)}  \\\hline
    引力  & pull + strength = \textit{gravitation}  & イン{\middel}リョク \mbox{(IN{\middel}RYOKU)}  \\\hline
    万引き  & ten thousand + pull = \textit{shoplifter}  & マン　び{\middel}き \mbox{(MAN bi{\middel}ki)}  \\\hline
    \notyet{引き受ける}  & pull + receive = \textit{to undertake}  & ひ{\middel}き　う{\middel}ける \mbox{(hi{\middel}ki u{\middel}keru)}  \\\hline
    \notyet{引き上げる}  & pull + upper = \textit{to pull up}  & ひ{\middel}き　あ{\middel}げる \mbox{(hi{\middel}ki a{\middel}geru)}  \\\hline
\end{somme}

\begin{sentence}
    \underline{やっと}{\;}\underline{\ruby{牛}{うし}}\ruby{を}{お}\underline{\ruby{川}{かわ}}{\;}\underline{から}{\;}\underline{\ruby{引}{ひ}き\ruby{上}{あ}げました}。\\\hline
yatto ushi o kawa kara hi{\middel}ki a{\middel}gemashita.\\\hline
    (finally) (cow) (river) (from) (pulled up)\\\hline
    =\textit{(We) finally pulled the cow from the river.}\\\hline
\end{sentence}
\end{multicols}
\vspace*{\fill}
\pagebreak
%%--------------------------------------------------------------------
%7--------------------------------------------------------------------
\vspace*{\fill}
\begin{multicols}{2}
\begin{glyph}{Second (秒)}{second}\label{glyph:second}
    Second (of time or arc).
\end{glyph}

\begin{comps}
\comprtwocone & 禾 + 少 \\\hline
    \comprthreecone & Two-branch tree + Few (\ref{glyph:few}) \\\hline
\end{comps}

\begin{remglyph}
    \hooglig{Seconds} until a \remcom{few} \remcom{crops} remained.\\\hline
\end{remglyph}

\begin{prons}
    \pronsrtwocone & \textbf{ビョウ (BY\={O})} &---\\\hline
\pronsrthreecone & --- &  ---\\\hline
\end{prons}

\begin{remprons}
    \hooglig{Seconds} untill it will \comon{be yours}.\\\hline
\end{remprons}

%{\middel}
% &  ()\\\hline
\begin{somme}
    一秒  & one + second = \textit{one second}  & イチ{\middel}ビョウ \mbox{(ICHI{\middel}BY\={O})}  \\\hline
    五秒  & five + second = \textit{five seconds}  & ゴ{\middel}ビョウ \mbox{(GO{\middel}BY\={O})}  \\\hline
    秒読み  & second + read = \textit{countdown}  & ビョウ　よ{\middel}み \mbox{(BY\={O} yo{\middel}mi)}  \\\hline
\end{somme}

\begin{sentence}
Note that the sample sentence contains one of the less common
    on-yomi for ``分'' (Entry 124) in the compound ``十五分''. \\\hline
    \underline{\ruby{十五分}{ジュウゴフン}}{\;}\underline{\ruby{二十秒}{ニジュウビョウ}}{\;}\underline{かかりました}。\\\hline
    J\={U}{\middel}GO{\middel}FUN NI{\middel}J\={U}{\middel}BY\={O} kakarimashita.\\\hline
    (fifteen minutes) (twenty seconds) (took)\\\hline
    =\textit{It took fifteen minutes and twenty seconds.}\\\hline
\end{sentence}
\end{multicols}
\vhrulefill{5pt}
%%--------------------------------------------------------------------
%8--------------------------------------------------------------------
\begin{multicols}{2}
\begin{glyph}{Shell (甲)}{shell}\label{glyph:shell}
Shell.\\\hline
This character rarely appears on its own, but proves valuable as a component.
\end{glyph}

\begin{comps}
\comprtwocone & 田 \\\hline
    \comprthreecone & Rice field (\ref{glyph:rice_field}) \\\hline
\end{comps}

\begin{remglyph}
---\\\hline
\end{remglyph}

\begin{prons}
    \pronsrtwocone & \textbf{コウ (K\={O})} & ---\\\hline
    \pronsrthreecone & カン (KAN) & --- \\\hline
\end{prons}

\begin{remprons}
    \hooglig{Shells} \comon{cause} \ucomon{cunning} schemes.\\\hline
\end{remprons}

%{\middel}
% &  ()\\\hline
\begin{somme}
    手の甲  & hand + shell = \textit{back of the hand}  & て{\middel}の{\middel}コウ \mbox{(te{\middel}no{\middel}K\={O})}  \\\hline
\end{somme}

\begin{sentence}
    \underline{\ruby{上田}{うえだ}さん}は\underline{\ruby{手}{て}の\ruby{甲}{コウ}}が\underline{かゆい}。\\\hline
    Ue{\middel}da-san wa te{\middel}no{\middel}K\={O} ga kayui.\\\hline
    (Ueda-san) (back of the hand) (itchy)\\\hline
    =\textit{The back of Ueda-san's hand is itchy.}\\\hline
\end{sentence}
\end{multicols}
\vspace*{\fill}
\pagebreak
%%--------------------------------------------------------------------
%9--------------------------------------------------------------------
\vspace*{\fill}
\begin{multicols}{2}
\begin{glyph}{Spirit (気)}{spirit}\label{glyph:spirit}
``Spirit'', ``air'', ``feeling'', and ``mood''.\\\hline
This kanji is found in many common words and expressions.
\end{glyph}

\begin{comps}
\comprtwocone & 气 + メ \\\hline
\comprthreecone & Steam/Air + ME (katakana) \\\hline
\end{comps}

\begin{remglyph}
\hooglig{Spiritual} \remcom{men} in the \remcom{steam} room.\\\hline
\end{remglyph}

\begin{prons}
\pronsrtwocone & \textbf{キ (KI)} & ---\\\hline
    \pronsrthreecone &  ケ (KE) &  --- \\\hline
\end{prons}

\begin{remprons}
    A \hooglig{spirit}-\comon{killing} \ucomon{keg}.\\\hline
\end{remprons}

%{\middel}
% &  ()\\\hline
\begin{somme}
    気  & \textit{spirit}  & キ \mbox{(KI)}  \\\hline
    元気  & basis + spirit = \textit{high spirits}  & ゲン{\middel}キ \mbox{(GEN{\middel}KI)}  \\\hline
    人気  & person + spirit = \textit{popularity}  & ニン{\middel}キ \mbox{(NIN{\middel}KI)}  \\\hline
    電気  & electric + spirit = \textit{electricity}  & デン{\middel}キ \mbox{(DEN{\middel}KI)}  \\\hline
    天気  & heaven + spirit = \textit{the weather}  & テン{\middel}キ \mbox{(TEN{\middel}KI)}  \\\hline
    気持ち  & spirit + hold = \textit{feeling}  & キ　も{\middel}ち \mbox{(KI mo{\middel}chi)}  \\\hline
\end{somme}

\begin{sentence}
The sample sentence is an extremely common greeting/question in
    Japanese.\\\hline
    \underline{\ruby{元気}{ゲンキ}}ですか。\\\hline
GEN{\middel}KI desu ka.\\\hline
    (high spirits)\\\hline
    =\textit{How are you?}\\\hline
\end{sentence}
\end{multicols}
\vhrulefill{5pt}
%%--------------------------------------------------------------------
%10-------------------------------------------------------------------
\begin{multicols}{2}
\begin{glyph}{Time (時)}{time}\label{glyph:time}
Time.\\\hline
This character will become familiar to you through its use in the hours of the
    day (the third and fourth compounds below are examples).
\end{glyph}

\begin{comps}
\comprtwocone & 日 + 寺  \\\hline
    \comprthreecone & Sun + Temple (\ref{glyph:temple}) \\\hline
\end{comps}

\begin{remglyph}
    \hooglig{Time} is calibrated in \remcom{temples} by means of the
    \remcom{sun}.\\\hline
\end{remglyph}

\begin{prons}
\pronsrtwocone & \textbf{ジ (JI)} & \textbf{とき (toki)}\\\hline
\pronsrthreecone & --- & --- \\\hline
\end{prons}

\begin{remprons}
    \hooglig{Time} the \comkun{torque} to \comon{jiggle}.\\\hline
\end{remprons}

%{\middel}
% &  ()\\\hline
\begin{somme}
    \rye{3}{Note how the final half of the second compound becomes voiced; this
    is an occasional feature of repeated compounds.}
    時  & \textit{time}  & とき \mbox{(toki)} \\\hline
    時々  & time + time = \textit{sometimes}  & とき{\middel}どき \mbox{(toki{\middel}doki)}  \\\hline
    四時  & four + time = \textit{four o' clock}  & よ{\middel}ジ \mbox{(yo{\middel}JI)} \\\hline
    九時半  & nine + time + half = \textit{nine-thirty}  & ク{\middel}ジ{\middel}ハン \mbox{(KU{\middel}JI{\middel}HAN)} \\\hline
    \notyet{同時}  & same + time = \textit{simultaneous}  &  ドウ{\middel}ジ \mbox{(D\={O}{\middel}JI)} \\\hline
\end{somme}

\begin{irrr}
    \notyet{時計}  & time + measure = \textit{watch; clock}  & と{\middel}ケイ \mbox{(to{\middel}KEI)}  \\\hline
\end{irrr}

\begin{sentence}
    \underline{\ruby{八時半}{ハチジハン}}に\underline{\ruby{東口}{ひがしぐち}}の\underline{\ruby{近}{ちか}く}で\underline{\ruby{会}{あ}いましょう}。\\\hline
    HACHI{\middel}JI{\middel}HAN ni higashi{\middel}guchi no chika{\middel}ku de a{\middel}imash\={o}.\\\hline
    (eight-thirty) (east entrance) (near) (let's meet)\\\hline
    =\textit{Let's meet near the east entrance at eight-thirty.}\\\hline
\end{sentence}
\end{multicols}
\vspace*{\fill}
\pagebreak
%%--------------------------------------------------------------------
%11-------------------------------------------------------------------
\vspace*{\fill}
\begin{multicols}{2}
\begin{glyph}{Speak (話)}{speak}\label{glyph:speak}
Speak, alon with the ideas of ``conversations'' and ``stories''.
\end{glyph}

\begin{comps}
\comprtwocone & 言 + 舌 \\\hline
    \comprthreecone & Say + Tongue \\\hline
\end{comps}

\begin{remglyph}
    \hooglig{``Speak''}, \remcom{says} the \remcom{tongue}.\\\hline
\end{remglyph}

\begin{prons}
\pronsrtwocone & \textbf{ワ (WA)} & \textbf{はなし、はな (hanashi, hana)}\\\hline
    \rye{3}{As in the final example below, the kun-yomi for this kanji is always
	voiced in the second position.}
\pronsrthreecone & --- & --- \\\hline
\end{prons}

\begin{remprons}
    \hooglig{Speak} how \comkun{hung-up} \comkun{Hanah shielded} us from being
    \comon{washed}.\\\hline
\end{remprons}

%{\middel}
% &  ()\\\hline
\begin{somme}
    話  & \textit{story}  & はなし \mbox{(hanashi)}  \\\hline
    話す  & \textit{to speak}  & はな{\middel}す \mbox{(hana{\middel}su)}  \\\hline
    話好き  & speak + like = \textit{talkative}  & はな{\middel}し　ず{\middel}き \mbox{(hana{\middel}shi zu{\middel}ki)}  \\\hline
    電話  & electric + speak = \textit{telephone}  &  デン{\middel}ワ \mbox{(DEN{\middel}WA)} \\\hline
    会話  & meet + speak = \textit{conversation}  & カイ{\middel}ワ \mbox{(KAI{\middel}WA)}  \\\hline
    立ち話  & stand + speak = \textit{chatting while standing}  & た{\middel}ち　ばなし \mbox{(ta{\middel}chi banashi)}  \\\hline
\end{somme}

\begin{sentence}
    \underline{\ruby{田中}{たなか}さん}は\underline{フランス\ruby{語}{ゴ}}\ruby{を}{お}\underline{\ruby{上手}{ジョウず}}に\underline{\ruby{話}{はな}せます}。\\\hline
    Ta{\middel}naka-san wa Furanso{\middel}GO o J\={O}{\middel}zu ni hana{\middel}semasu.\\\hline
    (Tanaka-san) (French) (well) (can speak)\\\hline
    =\textit{Tanaka-san can speak French well.}\\\hline
\end{sentence}
\end{multicols}
\vhrulefill{5pt}
%%--------------------------------------------------------------------
%12-------------------------------------------------------------------
\begin{multicols}{2}
\begin{glyph}{English (英)}{english}\label{glyph:english}
English.\\\hline
A secondary meaning relates to ``brilliance''.
\end{glyph}

\begin{comps}
\comprtwocone & 艸 + 央 \\\hline
    \comprthreecone & Grass + Center (\ref{glyph:center}) \\\hline
\end{comps}

\begin{remglyph}
    \hooglig{English} \remcom{grass} is \remcom{central}.\\\hline
\end{remglyph}

\begin{prons}
\pronsrtwocone & \textbf{エイ (EI)} & ---\\\hline
\pronsrthreecone & --- & --- \\\hline
\end{prons}

\begin{remprons}
    \hooglig{English} \comon{aim} for world domination.\\\hline
\end{remprons}

%{\middel}
% &  ()\\\hline
\begin{somme}
    英語  & English + words = \textit{English language}  & エイ{\middel}ゴ \mbox{(EI{\middel}GO)}  \\\hline
    英国  & English + country = \textit{England}  & エイ{\middel}コク \mbox{(EI{\middel}KOKU)}  \\\hline
    英会話  & English + meet + speak = \textit{English conversation}  & エイ{\middel}カイ{\middel}ワ \mbox{(EI{\middel}KAI{\middel}WA)}  \\\hline
    日英  & sun + English = \textit{Japan and England}  & ニチ{\middel}エイ \mbox{(NICHI{\middel}EI)}  \\\hline
\end{somme}

\begin{sentence}
    \underline{\ruby{英会話}{エイカイワ}}は\underline{\ruby{日本}{ニホン}}で\underline{\ruby{人気}{ニンキ}}があります。\\\hline
EI{\middel}KAI{\middel}WA wa NI{\middel}HON de NIN{\middel}KI ga arimasu.\\\hline
    (English conversation) (Japan) (popular) (is)\\\hline
    =\textit{English conversation is popular in Japan.}\\\hline
\end{sentence}
\end{multicols}
\vspace*{\fill}
\pagebreak
%%--------------------------------------------------------------------
%13-------------------------------------------------------------------
\vspace*{\fill}
\begin{multicols}{2}
\begin{glyph}{Instruction (訓)}{instruction}\label{glyph:instruction}
Instruction.\\\hline
This kanji tends to give a more serious or spiritual sense of ``teaching'' to
    the compounds in which it appears.
\end{glyph}

\begin{comps}
\comprtwocone & 言 + 川 \\\hline
    \comprthreecone & Say/Speech + River\\\hline
\end{comps}

\begin{remglyph}
    \hooglig{Instruct} the \remcom{river} to \remcom{say} something.\\\hline
\end{remglyph}

\begin{prons}
\pronsrtwocone & \textbf{クン (KUN)} & ---\\\hline
\pronsrthreecone & --- & --- \\\hline
\end{prons}

\begin{remprons}
    \hooglig{Instruct} the \comon{coon}?!\\\hline
\end{remprons}

%{\middel}
% &  ()\\\hline
\begin{somme}
    訓読み  & instruction + read = \textit{the kun-yomi}  & クン　よ{\middel}み \mbox{(KUN yo{\middel}mi)}  \\\hline
    訓話  & instruction + speak = \textit{moral discourse}  & クン{\middel}ワ \mbox{(KUN{\middel}WA)}  \\\hline
\end{somme}

\begin{sentence}
    \underline{これ}\ruby{を}{お}\underline{\ruby{訓読}{クンよ}み}で\underline{\ruby{読}{よ}んで}{\;}\underline{\ruby{下}{くだ}さい}。\\\hline
kore o KUN{\middel}yomi de yo{\middel}nde kuda{\middel}sai.\\\hline
    (this) (kun-yomi) (read) (please)\\\hline
    =\textit{Please read this with kun-yomi.}\\\hline
\end{sentence}
\end{multicols}
\vhrulefill{5pt}
%%--------------------------------------------------------------------
%14-------------------------------------------------------------------
\begin{multicols}{2}
\begin{glyph}{Push (押)}{push}\label{glyph:push}
``Push'', together with the ideas of ``suppression'' and ``restraint''.\\\hline
    This kanji also serves as the partner to 引 (Entry \ref{glyph:pull}) on door signs.
\end{glyph}

\begin{comps}
\comprtwocone & 手 + 田 \\\hline
    \comprthreecone & Hand (\ref{glyph:hand}) + Rice field
    (\ref{glyph:rice_field}) \\\hline
\end{comps}

\begin{remglyph}
    \hooglig{Push} your \remcom{hand} into the \remcom{rice fields}.\\\hline
\end{remglyph}

\begin{prons}
\pronsrtwocone & --- & \textbf{お (o)}\\\hline
    \pronsrthreecone & オウ (OU) & --- \\\hline
\end{prons}

\begin{remprons}
    \hooglig{Pushing} can be \ucomon{awefully} \comkun{obligatory}.\\\hline
\end{remprons}

%{\middel}
% &  ()\\\hline
\begin{somme}
    \rye{3}{As indicated by the last two examples below, this is another kanji
    that often shows up in the manner of 買.}
    押す  & \textit{to push}  & お{\middel}す \mbox{(o{\middel}su)}  \\\hline
    押し出す  & push + exit = \textit{to push out}  & お{\middel}し　だ{\middel}す \mbox{(o{\middel}shi da{\middel}su)}  \\\hline
    押し上げる  & push + upper = \textit{to push up}  & お{\middel}し　あ{\middel}げる \mbox{(o{\middel}shi a{\middel}geru)}  \\\hline
    押し開ける  & push + open = \textit{to push open}  & お{\middel}し　あ{\middel}ける \mbox{(o{\middel}shi a{\middel}keru)}  \\\hline
    押入れ  & push + enter = \textit{closet}  & おし　い{\middel}れ \mbox{(oshi i{\middel}re)}  \\\hline
    押売  & push + sell = \textit{high-pressure sales}  & おし{\middel}うり \mbox{(oshi{\middel}uri)}  \\\hline
\end{somme}

\begin{sentence}
    \underline{\ruby{押入}{おしい}れ}の\underline{\ruby{中}{なか}}に\underline{\ruby{物}{もの}}が\underline{\ruby{少}{すく}ない}。\\\hline
oshi i{\middel}re no naka ni mono ga suku{\middel}nai.\\\hline
    (closet) (middle) (things) (few)\\\hline
    =\textit{There aren't many things in the closet.}\\\hline
\end{sentence}
\end{multicols}
\vspace*{\fill}
\pagebreak
%%--------------------------------------------------------------------
%15-------------------------------------------------------------------
\vspace*{\fill}
\begin{multicols}{2}
\begin{glyph}{Year (年)}{year}\label{glyph:year}
Year.
\end{glyph}

\begin{comps}
\comprtwocone & 干 + 丿 + 一 \\\hline
    \comprthreecone & Dry + Slash + One \\\hline
\end{comps}

\begin{remglyph}
    \hooglig{Year} \remcom{one} \remcom{dries} out in a \remcom{slash}.\\\hline
\end{remglyph}

\begin{prons}
\pronsrtwocone & \textbf{ネン (NEN)} & \textbf{とし (toshi)}\\\hline
\pronsrthreecone & --- & --- \\\hline
\end{prons}

\begin{remprons}
    The \hooglig{year} the \comon{Nenets} acquired \comkun{Toshiba}.\\\hline
\end{remprons}

%{\middel}
% &  ()\\\hline
\begin{somme}
    年  & \textit{year}  & とし \mbox{(toshi)}  \\\hline
    年上  & year + upper = \textit{older}  & とし{\middel}うえ \mbox{(toshi{\middel}ue)}  \\\hline
    年下  & year + lower = \textit{younger}  & とし{\middel}した \mbox{(toshi{\middel}shita)}  \\\hline
    新年  & new + year = \textit{the New year}  & シン{\middel}ネン \mbox{(SHIN{\middel}NEN)}  \\\hline
    中年  & middle + year = \textit{middle age}  & チュウ{\middel}ネン \mbox{(CH\={U}{\middel}NEN)}  \\\hline
    一年生  & one + year + life = \textit{first-year student}  & イチ{\middel}ネン{\middel}セイ \mbox{(ICHI{\middel}NEN{\middel}SEI)}  \\\hline
\end{somme}

\begin{sentence}
    \underline{\ruby{青木}{あおき}さん}は\underline{\ruby{三年生}{ワンネンセイ}}です。\\\hline
Ao{\middel}ki-san wa SAN{\middel}NEN{\middel}SEI desu.\\\hline
    (Aoki-san) (third-year student)\\\hline
    =\textit{Aoki-san is a third-year student.}\\\hline
\end{sentence}
\end{multicols}
\vhrulefill{5pt}
%%--------------------------------------------------------------------
%16-------------------------------------------------------------------
\begin{multicols}{2}
\begin{glyph}{Come (来)}{come}\label{glyph:come}
Come.\\\hline
As the final four compounds below make clear, the related meanings of ``next''
    is also included.
\end{glyph}

\begin{comps}
\comprtwocone & 一 + 米 \\\hline
    \comprthreecone & One + Rice \\\hline
\end{comps}

\begin{remglyph}
    \hooglig{Come} \remcom{one} \remcom{rice} grain at a time.\\\hline
\end{remglyph}

\begin{prons}
\pronsrtwocone & \textbf{アイ (RAI)} & \textbf{く (ku)}\\\hline
    \pronsrthreecone & --- & き (ki) \\\hline
    \rye{3}{This character is infamous for being one of the only two irregular
    verbs in Japanese.  The character can be read variously as く, き, or こ,
    depending on which tense of the verb is being used.  The sample sentence
    offers an example of the polite past tense.  Note, however, an important
    exception in the second compound below (an extremely common word); when
    following the kanji 出, 来 is always read き.}
\end{prons}

\begin{remprons}
    \hooglig{Come} up with a \comkun{coup} to \ucomkun{keep} on \comon{rhyming}.\\\hline
\end{remprons}

%{\middel}
% &  ()\\\hline
\begin{somme}
    来る  & \textit{to come}  & く{\middel}る \mbox{(ku{\middel}ru)}  \\\hline
    出来る  & exit + come = \textit{to be able to}  & で　き{\middel}る \mbox{(de ki{\middel}ru)}  \\\hline
    来週  & come + week = \textit{next week}  & ライ{\middel}シュウ \mbox{(RAI{\middel}SH\={U})}  \\\hline
    来月  & come + moon = \textit{next month}  & ライ{\middel}ゲツ \mbox{(RAI{\middel}GETSU)}  \\\hline
    来年  & come + year = \textit{next year}  & ライ{\middel}ネン \mbox{(RAI{\middel}NEN)}  \\\hline
    来春  & come + spring = \textit{next spring}  & ライ{\middel}シュン \mbox{(RAI{\middel}SHUN)}  \\\hline
\end{somme}

\begin{sentence}
    \underline{\ruby{先生}{センセイ}}は\underline{\ruby{一週間}{イッシュウカン}}{\;}\underline{\ruby{前}{まえ}}に\underline{\ruby{来}{き}ました}。\\\hline
    SEN{\middel}SEI wa IS{\middel}SH\={U}{\middel}KAN mae ni ki{\middel}mashita.\\\hline
    (teacher) (one week) (before) (came)\\\hline
    =\textit{The teacher came a week ago.}\\\hline
\end{sentence}
\end{multicols}
\vspace*{\fill}
\pagebreak
%%--------------------------------------------------------------------
%17-------------------------------------------------------------------
\vspace*{\fill}
\begin{multicols}{2}
\begin{glyph}{Study (学)}{study}\label{glyph:study}
Study and learning.\\\hline
When used as a suffix, it is often equivalent to the English ``-ology''.  The
    final compound offers such an example.
\end{glyph}

\begin{comps}
\comprtwocone & 子 + 小 + 冖 \\\hline
    \comprthreecone & Child + Small + Cover/Crown \\\hline
\end{comps}

\begin{remglyph}
\hooglig{Study} of \remcom{small} \remcom{children's} \remcom{crowns}.\\\hline
\end{remglyph}

\begin{prons}
\pronsrtwocone & \textbf{ガク (GAKU)} & \textbf{まな (mana)}\\\hline
\pronsrthreecone & --- & --- \\\hline
\end{prons}

\begin{remprons}
    \hooglig{Study} how to \comkun{manage} all this \comon{guck}.\\\hline
\end{remprons}

%{\middel}
% &  ()\\\hline
\begin{somme}
    学ぶ  & \textit{to study}  & まな{\middel}ぶ \mbox{(mana{\middel}bu)}  \\\hline
    大学  & large + study = \textit{university}  & ダイ{\middel}ガク \mbox{(DAI{\middel}GAKU)}  \\\hline
    学生  & study + life = \textit{student}  & ガク{\middel}セイ \mbox{(GAKU{\middel}SEI)}  \\\hline
    休学  & rest + study = \textit{absence from school}  & キュウ{\middel}ガク \mbox{(KY\={U}{\middel}GAKU)}  \\\hline
    語学  & words + study = \textit{linguistics}  & ゴ{\middel}ガク \mbox{(GO{\middel}GAKU)}  \\\hline
    生物学  & life + thing + study = \textit{biology}  & セイ{\middel}ブツ{\middel}ガク \mbox{(SEI{\middel}BUTSU{\middel}GAKU)}  \\\hline
\end{somme}

\begin{sentence}
    \underline{この}{\;}\underline{\ruby{家}{いえ}}に\ruby{は}{わ}\underline{\ruby{大学生}{ダイガクセイ}}が\underline{\ruby{多}{おお}い}。\\\hline
    kono ie ni wa DAI{\middel}GAKU{\middel}SEI ga \={o}{\middel}i.\\\hline
    (this) (house) (university students) (many)\\\hline
    =\textit{There are many university students in this house.}\\\hline
\end{sentence}
\end{multicols}
\vhrulefill{5pt}
%%--------------------------------------------------------------------
%18-------------------------------------------------------------------
\begin{multicols}{2}
\begin{glyph}{Mind (意)}{mind}\label{glyph:mind}
Mind/Thought.\\\hline
This abstract character appears in many compounds conveying the ideas of
    ``will'', ``intention'', and ``opinion'', etc.
\end{glyph}

\begin{comps}
\comprtwocone & 音 + 心 \\\hline
    \comprthreecone & Sound + Heart \\\hline
\end{comps}

\begin{remglyph}
    \hooglig{Mind} and \remcom{heart} don't always \remcom{resonate}.\\\hline
\end{remglyph}

\begin{prons}
\pronsrtwocone & \textbf{イ (I)} & ---\\\hline
\pronsrthreecone & --- & --- \\\hline
\end{prons}

\begin{remprons}
    \hooglig{Mind} the \comon{indicators}.\\\hline
\end{remprons}

%{\middel}
% &  ()\\\hline
\begin{somme}
    意思  & mind + think = \textit{(one's) intent}  & イ{\middel}シ \mbox{(I{\middel}SHI)}  \\\hline
    意図  & mind + diagram = \textit{(one's) aim}  & イ{\middel}ト \mbox{(I{\middel}TO)}  \\\hline
    意外  & mind + outside = \textit{unforseen}  & イ{\middel}ガイ \mbox{(I{\middel}GAI)}  \\\hline
    意見  & mind + see = \textit{opinion}  & イ{\middel}ケン \mbox{(I{\middel}KEN)}  \\\hline
    意気  & mind + spirit = \textit{morale}  & イ{\middel}キ \mbox{(I{\middel}KI)}  \\\hline
    \notyet{注意}  & pour + mind = \textit{attention}  & チュウ{\middel}イ \mbox{(CH\={U}{\middel}I)}  \\\hline
\end{somme}

\begin{sentence}
    \underline{あなた}の\underline{\ruby{意見}{イケン}}が\underline{\ruby{分}{わ}かりません}。\\\hline
anata no I{\middel}KEN ga wa{\middel}karimasen.\\\hline
    (you) (opinion) (don't understand)\\\hline
    =\textit{I don't understand your opinion.}\\\hline
\end{sentence}
\end{multicols}
\vspace*{\fill}
\pagebreak
%%--------------------------------------------------------------------
%19-------------------------------------------------------------------
\vspace*{\fill}
\begin{multicols}{2}
\begin{glyph}{Again (又)}{again}\label{glyph:again}
Agains/Also.  A character that rarely appears on its own.\\\hline
The suggested immediate visual connection is ``ironing board''.
\end{glyph}

\begin{comps}
\comprtwocone & 又 \\\hline
    \comprthreecone & Again \\\hline
\end{comps}

\begin{remglyph}
Part of the Kanji radicals (\#{29}).\\\hline
\end{remglyph}

\begin{prons}
\pronsrtwocone & --- & \textbf{また (mata)}\\\hline
\pronsrthreecone & --- & --- \\\hline
\end{prons}

\begin{remprons}
    \hooglig{Again} and again the \comkun{matador} irked the bull.\\\hline
\end{remprons}

%{\middel}
% &  ()\\\hline
\begin{somme}
    \rye{3}{Note how the second kanji in the final example below is voiced.}
    又  & \textit{again}  & また \mbox{(mata)}  \\\hline
    又々  & again + again = \textit{once again}  & また{\middel}また \mbox{(mata{\middel}mata)}  \\\hline
    又聞き  & again + hear = \textit{hearsay}  & また{\middel}ぎき \mbox{(mata{\middel}giki)}  \\\hline
\end{somme}

\begin{sentence}
    The pronunciation of 来なかった in the sample sentence, incidentially, is
    the informal past tense reading of 来.\\\hline
    \underline{\ruby{北川}{きたがわ}さん}は\underline{\ruby{又}{また}}{\;}\underline{\ruby{来}{こ}なかった}よ。\\\hline
Kita{\middel}gawa-san wa mata ko{\middel}nakatta yo.\\\hline
    (Kitagawa-san) (again) (didn't come)\\\hline
    \textit{Kitagawa-san didn't come again!}\\\hline
\end{sentence}
\end{multicols}
\vhrulefill{5pt}
%%--------------------------------------------------------------------
%20-------------------------------------------------------------------
\begin{multicols}{2}
\begin{glyph}{Red (赤)}{red}\label{glyph:red}
Red.
\end{glyph}

\begin{comps}
\comprtwocone & 赤 \\\hline
    \comprthreecone & Red \\\hline
\end{comps}

\begin{remglyph}
Part of the Kanji radicals (\#{155}).\\\hline
\end{remglyph}

\begin{prons}
\pronsrtwocone & \textbf{セキ (SEKI)} & \textbf{あか (aka)}\\\hline
    \pronsrthreecone & シャク (SHAKU) & --- \\\hline
\end{prons}

\begin{remprons}
    \hooglig{Red} is \comkun{also known as} \comon{sexy}.  \ucomon{Shucks}!\\\hline
\end{remprons}

%{\middel}
% &  ()\\\hline
\begin{somme}
    赤  & \textit{red (noun)}  & あか \mbox{(aka)}  \\\hline
    赤い  & \textit{red (adjective)}  & あか{\middel}い \mbox{(aka{\middel}i)}  \\\hline
    赤ちゃん  & \textit{baby}  & あか{\middel}ちゃん \mbox{(aka{\middel}chan)}  \\\hline
    赤道  & red + road = \textit{equator}  & セキ{\middel}ドウ \mbox{(SEKI{\middel}D\={O})}  \\\hline
\end{somme}

\begin{sentence}
    \underline{あの}{\;}\underline{\ruby{赤}{あか}い}{\;}\underline{\ruby{花}{はな}}\ruby{は}{わ}\underline{\ruby{美}{うつく}しい}{\;}\underline{ですね}。\\\hline
    ano aka{\middel}i hana wa utsuku{\middel}shii desu ne.\\\hline
    (that) (red) (flower) (beautiful) (isn't it)\\\hline
    =\textit{That red flower is beautiful, isn't it?}\\\hline
\end{sentence}
\end{multicols}
\vhrulefill{5pt}
%%--------------------------------------------------------------------
%%--------------------------------------------------------------------
\vspace*{\fill}
\pagebreak
